\documentclass[12pt,a4paper]{article}

% ===========================================================================
% Paquetes
% ===========================================================================
\usepackage[utf8]{inputenc}
\usepackage[T1]{fontenc}
\usepackage[spanish,es-noshorthands]{babel}
\usepackage{geometry}
\geometry{top=2.5cm, bottom=2.5cm, left=2.8cm, right=2.8cm}

\usepackage{amsmath}
\usepackage{booktabs}
\usepackage{array}
\usepackage{xcolor}
\usepackage{listings}
\usepackage{float}
\usepackage{caption}
\usepackage{parskip}
\usepackage{microtype}
\usepackage{afterpage}   % para cambio de geometría en página específica
\usepackage{hyperref}
\hypersetup{
    colorlinks = true,
    linkcolor  = blue!60!black,
    urlcolor   = blue!60!black,
    pdfauthor  = {Equipo de Fundamentos de IA},
    pdftitle   = {Análisis de la FSM de la Máquina Expendedora}
}

% TikZ para diagramas UML de actividades
\usepackage{tikz}
\usetikzlibrary{shapes.geometric, arrows.meta, positioning, calc}

% Estilo de código C++
\lstdefinestyle{cpp}{
    language     = C++,
    basicstyle   = \ttfamily\footnotesize,
    keywordstyle = \color{blue!70!black}\bfseries,
    commentstyle = \color{green!50!black}\itshape,
    stringstyle  = \color{orange!80!black},
    numbers      = left,
    numberstyle  = \tiny\color{gray},
    numbersep    = 8pt,
    frame        = single,
    breaklines   = true,
    tabsize      = 4,
    captionpos   = b,
    aboveskip    = 8pt,
    belowskip    = 4pt
}


% ===========================================================================
\begin{document}
% ===========================================================================

% ---------------------------------------------------------------------------
% Portada
% ---------------------------------------------------------------------------
\begin{titlepage}
    \centering
    \vspace*{2cm}
    {\LARGE\bfseries Universidad Politécnica de Victoria}\\[0.5cm]
    {\large Ingeniería en Tecnologías de la Información}\\[1.5cm]
    \rule{\linewidth}{1.5pt}\\[0.6cm]
    {\huge\bfseries Análisis Técnico de la\\[0.3cm]
     Máquina de Estados Finitos (FSM)\\[0.3cm]
     de la Máquina Expendedora}\\[0.4cm]
    \rule{\linewidth}{1.5pt}\\[1.2cm]
    {\large\bfseries Materia:} Fundamentos de Inteligencia Artificial\\[0.3cm]
    {\large\bfseries Plataforma:} Arduino Mega 2560 (firmware C++, PlatformIO)\\[1.5cm]
    {\large Arturo Zapata \quad · \quad Fabián Ramírez \quad · \quad Equipo}\\[0.5cm]
    {\large\today}
    \vfill
\end{titlepage}

\tableofcontents
\newpage

% ===========================================================================
\section{Introducción y Contexto del Proyecto}
% ===========================================================================

La máquina expendedora es un sistema embebido autónomo implementado en un
\textbf{Arduino Mega 2560}. Todo el firmware corre en un único microcontrolador
sobre un \emph{super-loop} no bloqueante.

Los periféricos que gestiona el sistema son:
\begin{itemize}
    \item \textbf{Teclado matricial 4×4} (D22–D29) — entrada del usuario.
    \item \textbf{LCD 20×4 I$^2$C} (0x27) — retroalimentación visual.
    \item \textbf{Motores DC} vía PCA9685 (I$^2$C) — despacho de producto.
    \item \textbf{Lector RFID MFRC522} (SPI: D53, D8) — pago con tarjeta.
    \item \textbf{Barrera óptica de caída} (D2, pull-up) — confirmación de entrega del motor.
    \item \textbf{EEPROM interna} — persistencia de inventario y caja.
\end{itemize}

El núcleo del firmware es el módulo \texttt{VmFsm}, una Máquina de Estados
Finitos (FSM) con 13~estados que encapsula toda la lógica de negocio:
selección de producto, cobro, despacho y cálculo de cambio.

% ===========================================================================
\section{Punto de Entrada: \texttt{main.cpp}}
% ===========================================================================

\subsection{Arquitectura de Super-Loop}

El firmware sigue el patrón \emph{bare-metal super-loop}: \texttt{setup()}
inicializa los periféricos una sola vez y \texttt{loop()} se ejecuta
continuamente sin usar interrupciones ni RTOS.

\begin{lstlisting}[style=cpp, caption={Núcleo del super-loop en main.cpp}]
void setup() {
    Wire.begin();                  // I2C antes del LCD y PCA9685
    pinMode(VM_PIN_BARRIER, INPUT_PULLUP);
    display.begin();
    motor.begin();
    rfid.begin();
    fsm.begin();                   // carga EEPROM y entra al estado S0
}

void loop() {
    char key = keypad.getKey();    // sondeo no bloqueante del teclado
    if (key != NO_KEY)
        fsm.handleKey(key);        // despachar tecla a la FSM

    fsm.update();                  // avanzar la maquina de estados
}
\end{lstlisting}

\subsection{Inyección de Dependencias}

Todos los subsistemas se instancian como objetos estáticos globales y se
\emph{inyectan} al constructor de \texttt{VmFsm}:

\begin{lstlisting}[style=cpp, caption={Instanciación e inyección en main.cpp}]
static VmDisplay         display;
static VmEepromData      storage;
static VmMotorController motor;
static VmRfid            rfid;

static VmFsm fsm(storage, motor, rfid, onDisplay);
\end{lstlisting}

La función \texttt{onDisplay()} actúa como \emph{callback} de presentación:
la FSM nunca accede directamente al LCD; simplemente llama al puntero de
función \texttt{DisplayFn} que le fue entregado en construcción.

% ===========================================================================
\section{La Máquina de Estados: \texttt{VmFsm}}
% ===========================================================================

\subsection{Concepto y Rol}

\texttt{VmFsm} implementa el patrón \textbf{Mediador/Controlador}: centraliza
la lógica de negocio y coordina a los demás módulos sin que éstos se
conozcan entre sí. Es una \textbf{Máquina de Moore}: las salidas (pantallas,
motor, cambio) dependen únicamente del estado activo, y las transiciones
dependen de eventos (teclas, RFID, temporizadores, barrera óptica).

\subsection{Mecanismo de Transición}

\begin{enumerate}
    \item \texttt{enterState(FsmState next)} — cambia \texttt{\_state} y
          llama al gancho \texttt{onEnter*()} correspondiente.
    \item \texttt{update()} — llamado cada ciclo desde \texttt{loop()}.
          Sondea temporizadores, RFID y la barrera óptica para generar
          transiciones autónomas.
    \item \texttt{handleKey()} — traduce el carácter físico a una acción
          semántica y la delega al procesador del estado activo.
\end{enumerate}

\subsection{Temporización No Bloqueante}

Todos los temporizadores se implementan comparando \texttt{millis()} contra
una referencia guardada, sin usar \texttt{delay()}:

\begin{lstlisting}[style=cpp, caption={Patrón de temporizador no bloqueante}]
_motorTimer = millis();                           // arrancar

// En update():
if ((uint32_t)(millis() - _motorTimer) >= VM_TIMEOUT_MOTOR_MS)
    // ... accion de timeout
\end{lstlisting}

\begin{table}[H]
\centering
\begin{tabular}{lrp{5.5cm}}
\toprule
\textbf{Constante} & \textbf{Valor} & \textbf{Uso} \\
\midrule
\texttt{VM\_TIMEOUT\_INACTIVITY\_MS} & 180\,000 ms & Regreso a reposo en S3–S6 \\
\texttt{VM\_TIMEOUT\_MOTOR\_MS}      &  10\,000 ms & Tiempo máximo de giro del motor en S8 \\
\texttt{VM\_CAROUSEL\_INTERVAL\_MS}  &   2\,000 ms & Rotación de subpantalla del carrusel \\
\bottomrule
\end{tabular}
\caption{Temporizadores de la FSM}
\end{table}

\subsection{Interrelación con las Demás Clases}

\begin{table}[H]
\centering
\renewcommand{\arraystretch}{1.3}
\begin{tabular}{lp{9cm}}
\toprule
\textbf{Clase} & \textbf{Responsabilidad} \\
\midrule
\texttt{VmEepromData}      & Persistencia de slots (nombre, precio, stock) y caja de efectivo. \\
\texttt{VmMotorController} & Controla motores DC vía PCA9685; detecta entrega mediante barrera óptica (D2). \\
\texttt{VmRfid}            & Lectura de tarjetas MIFARE y escritura del saldo tras el cobro. \\
\texttt{VmKeypad}          & Interpreta caracteres del teclado 4×4 en acciones semánticas según el contexto. \\
\texttt{VmDisplay}         & Abstracción del LCD 20×4 I$^2$C; recibe 4 cadenas de texto. \\
\texttt{VmCarousel}        & Rota subpantallas cada 2\,s llamando al callback de display. \\
\texttt{VmChangeCalculator}& Algoritmo voraz de cálculo de cambio sobre las monedas de la caja. \\
\bottomrule
\end{tabular}
\caption{Módulos del firmware y sus responsabilidades}
\end{table}

% ===========================================================================
\section{Cálculo de Cambio: \texttt{VmChangeCalculator}}
% ===========================================================================

\subsection{Denominaciones Aceptadas}

El sistema maneja cuatro denominaciones de moneda mexicana (máximo \$10):

\begin{center}
\$1.00 (100c) \quad \$2.00 (200c) \quad \$5.00 (500c) \quad \$10.00 (1000c)
\end{center}

En el teclado, la correspondencia es: \texttt{1}=\$1, \texttt{2}=\$2,
\texttt{3}=\$5, \texttt{4}=\$10. Las teclas 5–9 se ignoran.

\subsection{Algoritmo Greedy Decreciente}

El cambio se calcula con un \textbf{algoritmo voraz}: se intentan primero
las monedas de mayor denominación, reduciendo el monto restante hasta
completarlo o agotar existencias.

\begin{lstlisting}[style=cpp, caption={Núcleo del algoritmo greedy en VmChangeCalculator::calculate()}]
const uint32_t denominaciones[4] = { 1000u, 500u, 200u, 100u };

for (uint8_t i = 0; i < 4 && total > 0u; i++) {
    uint32_t disponible = 0;
    cashBox.getCoinStock(denominaciones[i], disponible);

    uint32_t usar = 0;
    while (usar < disponible && total >= denominaciones[i]) {
        total -= denominaciones[i];
        usar++;
    }
    if (usar > 0u) {
        cashBox.deductCoins(denominaciones[i], usar);
        *contadores[i] = usar;
    }
}
if (total > 0u) return false;  // cambio no realizable
\end{lstlisting}

\subsection{Visualización del Cambio en Carrusel}

Tras un pago exitoso en efectivo, la FSM construye un \textbf{carrusel de
slides}: un slide inicial de éxito seguido de uno por cada denominación
no nula (de mayor a menor). Las denominaciones con cero monedas se omiten:

\begin{lstlisting}[style=cpp, caption={Construcción dinámica del carrusel de cambio en onEnterPantallaFin()}]
_carousel.addScreen(finSuccessTrampoline);   // slide 0: Compra Exitosa

const uint32_t denoms[4] = {1000u, 500u, 200u, 100u};
for (uint8_t i = 0; i < 4; i++) {
    if (qtys[i] > 0u) {
        _changeDenoms[_changeSlides] = denoms[i];
        _changeQtys[_changeSlides]   = qtys[i];
        _changeSlides++;
        _carousel.addScreen(changeScreenTrampoline);
    }
}
\end{lstlisting}

Ejemplo: si el cambio es \$7.00, el carrusel muestra tres slides:
\textit{Compra Exitosa} → \textit{\$5.00 × 1 pieza} → \textit{\$2.00 × 1 pieza}.

\subsection{Decisiones de Diseño}

\begin{itemize}
    \item \textbf{Aritmética entera en centavos}: sin errores de coma flotante en AVR de 8~bits.
    \item \textbf{Descuento inmediato en EEPROM}: cada \texttt{deductCoins()} persiste al instante.
    \item \textbf{Retorno \texttt{false} si no hay fondos}: mensaje 5~s y regreso a reposo.
\end{itemize}

% ===========================================================================
\section{Los 13 Estados de la FSM}
% ===========================================================================

\subsection{Tabla Resumen}

\begin{table}[H]
\centering
\small
\renewcommand{\arraystretch}{1.25}
\begin{tabular}{clp{4.2cm}p{4.2cm}}
\toprule
\textbf{ID} & \textbf{Nombre} & \textbf{Propósito} & \textbf{Transiciones de salida} \\
\midrule
$S_0$  & \texttt{S0\_ARRANQUE}      & Inicialización; muestra mensaje de bienvenida. & → $S_2$ inmediatamente. \\
$S_1$  & \texttt{S1\_FALLA\_INTERNA}& Error crítico de EEPROM; sistema detenido.     & Estado terminal. \\
$S_2$  & \texttt{S2\_REPOSO}        & Carrusel del catálogo (2 slides de 2 productos). & → $S_3$ (tecla 1–4). \\
$S_3$  & \texttt{S3\_SEL\_CANAL}    & Confirmar producto elegido.  & → $S_4$ [A] / $S_2$ [*] / timeout. \\
$S_4$  & \texttt{S4\_SEL\_PAGO}     & Elegir método: efectivo o RFID. & → $S_5$ [A] / $S_6$ [B] / $S_2$ [*]. \\
$S_5$  & \texttt{S5\_ESP\_EFECTIVO} & Acumular monedas (teclas 1–4: \$1/\$2/\$5/\$10). & → $S_7$ (monto $\geq$ precio) / $S_2$ [B] / timeout. \\
$S_6$  & \texttt{S6\_ESP\_RFID}     & Sondear RFID y cobrar saldo. & → $S_7$ (OK) / $S_2$ (error/timeout/[B]). \\
$S_7$  & \texttt{S7\_RESERVADA}     & Reservar stock y arrancar motor. & → $S_8$ (OK) / $S_{10}$ (motor falló) / $S_2$ (sin stock, 1.5\,s). \\
$S_8$  & \texttt{S8\_DISPENSANDO}   & Motor activo; esperar barrera óptica.  & → $S_9$ (barrera OK) / $S_{10}$ (timeout 10\,s). \\
$S_9$  & \texttt{S9\_CONFIRMADA}    & Entrega confirmada; esperar 1\,s. & → $S_{11}$ (efectivo) / $S_{12}$ (RFID). \\
$S_{10}$& \texttt{S10\_FALLA\_DISP} & Rollback de stock; mensaje de error 3\,s. & → $S_2$. \\
$S_{11}$& \texttt{S11\_CALC\_CAMBIO}& Algoritmo greedy; preparar slides de cambio. & → $S_{12}$ (éxito) / $S_2$ (sin fondos, 5\,s). \\
$S_{12}$& \texttt{S12\_PANTALLA\_FIN}& Carrusel: éxito + una diapositiva por denominación no nula. & → $S_2$. \\
\bottomrule
\end{tabular}
\caption{Los 13 estados de la FSM de la máquina expendedora}
\end{table}

\subsection{Descripción Detallada de Cada Estado}

\paragraph{$S_0$ — Arranque.} Muestra mensaje de bienvenida; en el primer
ciclo de \texttt{update()} transita a $S_2$.

\paragraph{$S_1$ — Falla Interna.} Estado terminal por error de EEPROM.
Requiere reinicio físico del sistema.

\paragraph{$S_2$ — Reposo.} Carrusel con dos slides: el primero muestra
los productos 1 y 2; el segundo, los productos 3 y 4 (cada 2\,s).

\paragraph{$S_3$ — Selección de Canal.} Consulta \texttt{getSlot()} para
verificar que el slot está habilitado. Muestra nombre y precio del producto.

\paragraph{$S_4$ — Selección de Pago.} El usuario elige la pasarela de cobro.
Si RFID no está disponible, vuelve a esta pantalla con un aviso.

\paragraph{$S_5$ — Espera de Efectivo.} Teclas 1–4 representan monedas de
\$1, \$2, \$5 y \$10. Cada pulsación llama a \texttt{addCoins()} y actualiza
el acumulado. Al cubrir el precio transita automáticamente a $S_7$.
La pantalla muestra lo insertado, lo que falta y las teclas válidas.

\paragraph{$S_6$ — Espera de RFID.} \texttt{VmRfid::poll()} se llama en cada
ciclo sin bloquear. Saldo suficiente → \texttt{deductBalance()} → $S_7$.
Saldo insuficiente o error de lectura → mensaje y $S_2$.

\paragraph{$S_7$ — Reservada.} Punto de no retorno transaccional:
(1) \texttt{reserveStock()} decrementa el inventario en EEPROM;
(2) \texttt{motor.start()} verifica que la barrera óptica esté libre y enciende el motor.
Si \texttt{start()} falla, transita a $S_{10}$.

\paragraph{$S_8$ — Dispensando.} Supervisa \texttt{motor.isBusy()} en cada
ciclo. La barrera óptica al detectar la caída del producto detiene el motor
y reporta \texttt{VM\_RESULT\_DELIVERED}. Un timeout de 10\,s protege contra atascos.

\paragraph{$S_9$ — Confirmada.} Muestra pantalla de éxito 1\,s. Luego,
si el pago fue en efectivo, avanza a $S_{11}$; si fue RFID (sin cambio), a $S_{12}$.

\paragraph{$S_{10}$ — Falla de Despacho.} Si \texttt{\_stockReserved} es
\texttt{true}, llama a \texttt{releaseStock()} para devolver la unidad al
inventario. Muestra mensaje de error 3\,s y retorna a $S_2$.

\paragraph{$S_{11}$ — Cálculo de Cambio.} Invoca
\texttt{VmChangeCalculator::calculate()}. En caso de éxito, construye los
arrays \texttt{\_changeDenoms[]} y \texttt{\_changeQtys[]} para el carrusel.
Si no hay fondos suficientes, muestra aviso y espera 5\,s antes de volver a $S_2$.

\paragraph{$S_{12}$ — Pantalla Final.} El carrusel muestra: slide de éxito,
seguido de un slide por cada denominación no nula (de \$10 a \$1) o el
resumen de cobro RFID. Regresa a $S_2$ al completar todos los slides + 1\,s.

% ===========================================================================
\section{Diagrama de Estados UML}
% ===========================================================================

Cada estado muestra sus condiciones de entrada y de salida con el estado destino.
La Figura~\ref{fig:estadosuml} presenta el diagrama completo en la siguiente pagina.

\clearpage
\pdfpagewidth=54cm
\pdfpageheight=30cm
\paperwidth=54cm
\paperheight=30cm
\textwidth=52cm
\textheight=28cm
\vsize=28cm
\hsize=52cm
\pdfhorigin=1cm
\pdfvorigin=1cm
\oddsidemargin=0cm
\topmargin=0cm
\headheight=0cm
\headsep=0cm
\thispagestyle{empty}

\begin{center}
\tikzset{
  sbox/.style={
    rectangle, rounded corners=4pt,
    draw=blue!55!black, thick,
    fill=blue!6,
    text width=3.8cm,
    minimum height=3.0cm,
    align=center,
    font=\fontsize{6.8}{8.5}\selectfont\sffamily,
    inner sep=4pt
  },
  failbox/.style={
    sbox, fill=red!6, draw=red!55!black
  },
  startcirc/.style={
    circle, fill=black, minimum size=11pt, inner sep=0
  },
  arr/.style={
    ->, >=Stealth, draw=gray!70!black, line width=0.7pt
  },
  lbl/.style={
    font=\fontsize{6.2}{7.5}\selectfont\sffamily,
    fill=white, inner sep=1.8pt, align=center
  }
}

\begin{tikzpicture}[font=\fontsize{6.8}{8.5}\selectfont\sffamily]

% Inicio
\node[startcirc] (start) at (-2.4, 0) {};

% Fila principal y=0
\node[sbox] (s0) at (0, 0) {
  \textbf{S0: ARRANQUE}\\[2pt]
  \rule{\linewidth}{0.3pt}\\[1pt]
  \textit{Entra:} encendido del\\
  microcontrolador\\[2pt]
  \rule{\linewidth}{0.3pt}\\[1pt]
  \textit{Sale:}\\
  EEPROM OK $\to$ S2\\
  EEPROM falla $\to$ S1
};

\node[sbox] (s2) at (5.8, 0) {
  \textbf{S2: REPOSO}\\[2pt]
  \rule{\linewidth}{0.3pt}\\[1pt]
  \textit{Entra:} arranque / cancel. /\\
  timeout / error / fin compra\\[2pt]
  \rule{\linewidth}{0.3pt}\\[1pt]
  \textit{Sale:}\\
  tecla 1--4 $\to$ S3
};

\node[sbox] (s3) at (11.6, 0) {
  \textbf{S3: SEL\_CANAL}\\[2pt]
  \rule{\linewidth}{0.3pt}\\[1pt]
  \textit{Entra:} tecla 1--4 en S2\\[2pt]
  \rule{\linewidth}{0.3pt}\\[1pt]
  \textit{Sale:}\\
  {[A]} confirmar $\to$ S4\\
  {[*]} / timeout $\to$ S2
};

\node[sbox] (s4) at (17.4, 0) {
  \textbf{S4: SEL\_PAGO}\\[2pt]
  \rule{\linewidth}{0.3pt}\\[1pt]
  \textit{Entra:} {[A]} en S3\\[2pt]
  \rule{\linewidth}{0.3pt}\\[1pt]
  \textit{Sale:}\\
  {[A]} efectivo $\to$ S5\\
  {[B]} RFID $\to$ S6\\
  {[*]} $\to$ S2
};

\node[sbox] (s7) at (25.6, 0) {
  \textbf{S7: RESERVADA}\\[2pt]
  \rule{\linewidth}{0.3pt}\\[1pt]
  \textit{Entra:} monto cubierto\\
  (efectivo o RFID OK)\\[2pt]
  \rule{\linewidth}{0.3pt}\\[1pt]
  \textit{Sale:}\\
  stock OK + motor $\to$ S8\\
  sin stock (1.5s) $\to$ S2\\
  motor falla $\to$ S10
};

\node[sbox] (s8) at (31.4, 0) {
  \textbf{S8: DISPENSANDO}\\[2pt]
  \rule{\linewidth}{0.3pt}\\[1pt]
  \textit{Entra:} motor.start()\\
  exitoso\\[2pt]
  \rule{\linewidth}{0.3pt}\\[1pt]
  \textit{Sale:}\\
  barrera optica $\to$ S9\\
  timeout 10s $\to$ S10
};

\node[sbox] (s9) at (37.2, 0) {
  \textbf{S9: CONFIRMADA}\\[2pt]
  \rule{\linewidth}{0.3pt}\\[1pt]
  \textit{Entra:} barrera optica\\
  confirma caida\\[2pt]
  \rule{\linewidth}{0.3pt}\\[1pt]
  \textit{Sale:}\\
  1s + efectivo $\to$ S11\\
  1s + RFID $\to$ S12
};

\node[sbox] (s12) at (43.0, 0) {
  \textbf{S12: PANTALLA\_FIN}\\[2pt]
  \rule{\linewidth}{0.3pt}\\[1pt]
  \textit{Entra:} RFID confirmado /\\
  cambio calculado (S11)\\[2pt]
  \rule{\linewidth}{0.3pt}\\[1pt]
  \textit{Sale:}\\
  fin carrusel $\to$ S2
};

% Fila superior: S5 (y = +5.2)
\node[sbox] (s5) at (21.5, 5.2) {
  \textbf{S5: ESP\_EFECTIVO}\\[2pt]
  \rule{\linewidth}{0.3pt}\\[1pt]
  \textit{Entra:} {[A]} en S4\\
  teclas 1--4: \$1/\$2/\$5/\$10\\[2pt]
  \rule{\linewidth}{0.3pt}\\[1pt]
  \textit{Sale:}\\
  monto $\ge$ precio $\to$ S7\\
  {[B]} / timeout $\to$ S2
};

% Fila inferior (y = -5.2)
\node[failbox] (s1) at (0, -5.2) {
  \textbf{S1: FALLA\_INTERNA}\\[2pt]
  \rule{\linewidth}{0.3pt}\\[1pt]
  \textit{Entra:} EEPROM.begin()\\
  falla (retorna false)\\[2pt]
  \rule{\linewidth}{0.3pt}\\[1pt]
  \textit{Sale:} ninguna\\
  (bloqueo del sistema)
};

\node[sbox] (s6) at (21.5, -5.2) {
  \textbf{S6: ESP\_RFID}\\[2pt]
  \rule{\linewidth}{0.3pt}\\[1pt]
  \textit{Entra:} {[B]} en S4\\[2pt]
  \rule{\linewidth}{0.3pt}\\[1pt]
  \textit{Sale:}\\
  RFID OK + saldo $\ge$ precio $\to$ S7\\
  error / {[B]} / timeout $\to$ S2
};

\node[sbox] (s10) at (29.5, -5.2) {
  \textbf{S10: FALLA\_DISP}\\[2pt]
  \rule{\linewidth}{0.3pt}\\[1pt]
  \textit{Entra:} motor falla (S7)\\
  o timeout 10s (S8)\\[2pt]
  \rule{\linewidth}{0.3pt}\\[1pt]
  \textit{Sale:}\\
  rollback stock + 3s $\to$ S2
};

\node[sbox] (s11) at (38.0, -5.2) {
  \textbf{S11: CALC\_CAMBIO}\\[2pt]
  \rule{\linewidth}{0.3pt}\\[1pt]
  \textit{Entra:} 1s en S9\\
  (pago en efectivo)\\[2pt]
  \rule{\linewidth}{0.3pt}\\[1pt]
  \textit{Sale:}\\
  cambio OK $\to$ S12\\
  sin fondos (5s) $\to$ S2
};

% === TRANSICIONES ORTOGONALES ===

% start -> S0
\draw[arr] (start) -- (s0.west);

% S0 -> S2
\draw[arr] (s0.east) -- node[lbl, above] {EEPROM OK} (s2.west);

% S0 -> S1
\draw[arr] (s0.south) -- node[lbl, right] {EEPROM falla} (s1.north);

% S2 -> S3
\draw[arr] (s2.east) -- node[lbl, above] {tecla 1--4} (s3.west);

% S3 -> S4
\draw[arr] (s3.east) -- node[lbl, above] {[A] confirmar} (s4.west);

% S3 -> S2 (cancel/timeout)
\draw[arr] (s3.north) -- ++(0, 1.0) -| node[lbl, above, pos=0.25] {[*] / timeout} ($(s2.north)+(1.2,0)$);

% S4 -> S2 (cancel)
\draw[arr] ($(s4.north)+(-0.8,0)$) -- ++(0, 1.8) -| node[lbl, above, pos=0.25] {[*] cancelar} ($(s2.north)+(0.6,0)$);

% S4 -> S5 (efectivo)
\draw[arr] ($(s4.north)+(0.8,0)$) |- node[lbl, left, pos=0.6] {[A] efectivo} (s5.west);

% S4 -> S6 (RFID)
\draw[arr] ($(s4.south)+(0.8,0)$) |- node[lbl, left, pos=0.6] {[B] RFID} (s6.west);

% S5 -> S7 (monto cubierto)
\draw[arr] (s5.east) -| node[lbl, right, pos=0.3] {monto $\ge$ precio} ($(s7.north)+(-0.8,0)$);

% S5 -> S2 (cancel/timeout) - ruta alta sobre S5
\draw[arr] (s5.north) -- ++(0, 1.0) -| node[lbl, above, pos=0.2] {[B] / timeout} ($(s2.north)+(0.0,0)$);

% S6 -> S7 (RFID OK)
\draw[arr] (s6.east) -| node[lbl, right, pos=0.3] {RFID OK + saldo OK} ($(s7.south)+(-0.8,0)$);

% S6 -> S2 (error/cancel)
\draw[arr] (s6.south) -- ++(0, -1.8) -| node[lbl, below, pos=0.2] {error / [B] / timeout} ($(s2.south)+(0.8,0)$);

% S7 -> S8
\draw[arr] (s7.east) -- node[lbl, above] {stock OK + motor} (s8.west);

% S7 -> S2 (sin stock) - pasa limpiamente por arriba de S5
\draw[arr] ($(s7.north)+(0.5,0)$) -- ++(0, 6.9) -| node[lbl, above, pos=0.15] {sin stock (1.5s)} ($(s2.north)+(-0.6,0)$);

% S7 -> S10 (motor falla)
\draw[arr] ($(s7.south)+(0.5,0)$) |- node[lbl, above, pos=0.7] {motor falla} ($(s10.north)+(-0.7,0)$);

% S8 -> S9
\draw[arr] (s8.east) -- node[lbl, above] {barrera OK} (s9.west);

% S8 -> S10 (timeout 10s)
\draw[arr] (s8.south) |- node[lbl, above, pos=0.7] {timeout 10s} ($(s10.north)+(0.7,0)$);

% S9 -> S12 (RFID)
\draw[arr] (s9.east) -- node[lbl, above] {1s + RFID} (s12.west);

% S9 -> S11 (efectivo)
\draw[arr] (s9.south) -- node[lbl, right] {1s + efectivo} (s11.north);

% S10 -> S2 (rollback)
\draw[arr] (s10.south) -- ++(0, -1.0) -| node[lbl, below, pos=0.15] {rollback + 3s} ($(s2.south)+(0,0)$);

% S11 -> S12 (cambio OK)
\draw[arr] (s11.east) -| node[lbl, right, pos=0.4] {cambio OK} (s12.south);

% S11 -> S2 (sin fondos)
\draw[arr] (s11.south) -- ++(0, -2.6) -| node[lbl, below, pos=0.12] {sin fondos (5s)} ($(s2.south)+(-0.8,0)$);

% S12 -> S2 (fin carrusel) - ruta superior sin cruces
\draw[arr] (s12.north) -- ++(0, 7.6) -| node[lbl, above, pos=0.12] {fin carrusel $\to$ S2} ($(s2.north)+(-1.2,0)$);

\end{tikzpicture}
\vspace{0.4cm}

\captionof{figure}{Diagrama de estados UML de la FSM de la máquina expendedora (13 estados). Cada estado muestra sus condiciones de entrada (\textit{Entra:}) y de salida (\textit{Sale:}) con el estado destino. Todas las transiciones siguen trayectorias ortogonales.}
\label{fig:estadosuml}
\end{center}

\clearpage
\pdfpagewidth=21cm
\pdfpageheight=29.7cm
\paperwidth=21cm
\paperheight=29.7cm
\vsize=24.7cm
\hsize=15.4cm
\restoregeometry


% ===========================================================================
\section{Conclusiones}
% ===========================================================================

El diseño de la FSM cumple los objetivos de un sistema embebido robusto
y legible:

\begin{enumerate}
    \item \textbf{Legibilidad lineal}: los 13 estados están numerados
          $S_0$–$S_{12}$ y el código fuente sigue el mismo orden de arriba
          a abajo, facilitando la lectura y el mantenimiento.
    \item \textbf{Sin sobreingeniería}: cada clase tiene una responsabilidad
          única. Se eliminaron el registro transaccional (EEPROM) y el sensor
          de puerta, que no corresponden al hardware real actual.
    \item \textbf{Visualización correcta del cambio}: el carrusel de $S_{12}$
          muestra dinámicamente solo las denominaciones no nulas, de mayor
          a menor, una por slide, facilitando la lectura al usuario.
    \item \textbf{Denominaciones reales}: el sistema acepta exclusivamente
          monedas de \$1, \$2, \$5 y \$10 (teclas 1–4), coherente con el
          hardware de la caja real.
    \item \textbf{Temporización no bloqueante}: todos los delays se implementan
          con \texttt{millis()}, permitiendo que el super-loop permanezca
          siempre receptivo a eventos externos.
    \item \textbf{Carrusel corregido}: el módulo \texttt{VmCarousel} incrementa
          correctamente el índice en cada ciclo y muestra la primera slide
          de forma inmediata al entrar a un nuevo estado.
\end{enumerate}

\end{document}
